\documentclass[11pt]{article}

\usepackage[margin=1in]{geometry}
\usepackage[T1]{fontenc}
\usepackage[utf8]{inputenc}
\usepackage{lmodern}
\usepackage{microtype}
\usepackage{amsmath,amssymb,amsthm}
\usepackage{booktabs}
\usepackage{enumitem}
\usepackage{xurl}
\usepackage{listings}
\usepackage[hidelinks]{hyperref}

\lstset{basicstyle=\footnotesize\ttfamily,breaklines=true,frame=single,columns=fullflexible}

\newtheorem{definition}{Definition}
\newtheorem{remark}{Remark}

\title{ABOM and OINL for Anonymity Claims:\\
Proof-Carrying Odds\,--\,Inflation Budgets Without Publishing Full Artifacts}
\author{}
\date{February 28, 2026 (archive v133)}

\begin{document}
\maketitle

\begin{abstract}
This series publishes brief, referee-grade papers, but a credible anonymity claim typically depends on a large artifact surface: measurement scripts, traces, compiler knobs, certificates, and verifier outputs.
Publicly shipping \emph{all} artifacts is often infeasible (space, privacy, operational constraints).

We propose two lightweight documentation objects that make claims checkable even in a source-only release.
An \emph{Anonymity Bill of Materials} (ABOM) is a signed, machine-checkable manifest that commits (by digest) to the artifacts and threat-model assumptions behind a claim.
An \emph{Odds-Inflation Nutrition Label} (OINL) is a human-readable summary extracted from the ABOM: it states the protected interface (ENF-ID), the per-epoch budget representation, and the intended window/threat declaration (TW-ID).
Together they support reproducible \emph{interfaces}: others can diff labels across releases, verify signatures and digests, and request the full artifact archive only when needed.
\end{abstract}

\paragraph{Scope.}
We isolate and formalize \emph{ABOM and OINL for Anonymity Claims} as a reusable contract surface in the anonymity / Anonymous-DHT series.

\paragraph{Imports.}
We treat odds inflation and endpoint sizing as imported (Anonymity~A; Synthesis~5).\cite{series:oddsinflation,series:endpointbridge}
We adopt the contract-object vocabulary (Synthesis~0) as the schema surface to which ABOM commits; CPPC manifests are a control-plane sibling of ABOM (AnonDHT~State~3).\cite{series:contractobjects,series:state3}
Interface declarations (exposure normal forms / ENF-ID) and threat/window declarations (TW-ID) are imported from the trace-interface and threat-window synthesis notes.\cite{series:traceifaces,series:threatwindows}
Replay-plan and state-declaration equivalence/drift rules are imported from Synthesis~18.\cite{series:planstateequiv}
For supply-chain style release receipts and transparency anchoring, see Release~A and W-Congestion-EQ.\cite{series:releaseproperty,series:wcongeq}

\paragraph{Artifact policy.} This archive is source-only; supporting artifacts (logs, traces, scripts, certificates, and verifier outputs) are available on request as a single bundle, committed by ABOM/OINL manifests (this paper) and governed by the series artifact policy (Synthesis~6).\cite{series:artifactpolicy}

\paragraph{Exports.}
Exports two schemas: (i) \emph{ABOM} (signed manifest of claim semantics and artifact digests) and
(ii) \emph{OINL} (a derived, human-readable label with claim id, projection, adjacency, budget representation, and horizon).


\section{Design goals}
ABOM borrows its ``bill of materials'' stance from software supply-chain practice: rather than publishing every artifact, we publish a signed manifest that commits to them and to their semantics.
This parallels SBOM standards (e.g., SPDX, CycloneDX) and attestation frameworks such as in-toto, but with anonymity-specific fields (threat windows, transcript projections, and budget accountants).\cite{ntiaSBOM2019,spdx301,cyclonedx,intoto2019}

ABOM/OINL are \emph{interfaces}, not proofs.
They aim to make three failure modes harder:
\begin{enumerate}[leftmargin=*]
\item \textbf{Silent drift:} knobs or evidence change without an auditable trail.
\item \textbf{Cherry-picked evidence:} a claim is supported by a subset of runs while other runs contradict it.
\item \textbf{Ambiguous semantics:} a budget is stated in a representation that downstream readers cannot interpret.
\end{enumerate}

\section{ABOM: the source of truth}
\begin{definition}[ABOM]\label{def:abom}
ABOM is a contract object in the shared vocabulary of Synthesis~0; this note records the anonymity-specific minimum surface.\cite{series:contractobjects}
An ABOM is a signed manifest that commits (by digest) to: a claim identifier and version; an ENF-ID naming the protected transcript projection; a TW-ID naming the threat/window declaration; an adjacency declaration; a budget representation; replay hooks (plan identifiers/spec digests); optional state-surface declarations when guarantees depend on long-lived control-plane state; artifact digests/availability flags; and a signature.
\end{definition}

\paragraph{Minimal schema sketch.}
The schema is intentionally permissive; it is a contract object for auditors.
A minimal ABOM (YAML/JSON) might contain:
\begin{lstlisting}
{
  "abom_version": "0.3",
  "claim_id": "anondht.destination_privacy.v1",
  "paper_ref": "anonymity_series/paperB_anonymous_dht_profiling",
  "exposure_nf_id": "sha256:<ENF-ID-hex>",
  "tw_id": "sha256:<TW-ID-hex>",
  "replay_hook": {"plan_id": "plan.replay.v1", "plan_spec_id": "sha256:<PLAN-SPEC-hex>"},
  "state_decl_id": "sha256:<STATE-DECL-hex>",
  "adjacency": "change-one-destination-class",
  "budget": {"type": "odds_inflation_bits", "epsilon_bits": 0.35, "delta": 1e-6, "unit": "lookup"},
  "artifacts": [
    {"name": "verifier_report.json", "sha256": "..."},
    {"name": "trace_bundle.tar.zst", "sha256": "...", "availability": "on_request"}
  ],
  "signing": {"scheme": "ed25519", "pubkey": "...", "signature": "..."}
}
\end{lstlisting}

\begin{remark}[Digests are the point]
ABOM allows a publication to be ``source-only'' while still being checkable: if an author later provides the private artifact archive, third parties can verify that the delivered artifacts match the published digests.
\end{remark}

\section{OINL: a label derived from ABOM}
OINL is a compact table extracted from the ABOM.
It is the object we expect most readers to consume and compare.

\begin{table}[t]
\centering
\begin{tabular}{@{}ll@{}}
\toprule
\textbf{OINL field} & \textbf{Meaning}\\
\midrule
Claim id & Stable string used for diffing and change control \\
Exposure (ENF-ID) & \texttt{exposure\_nf\_id} naming what part of the transcript is budgeted \\
Adjacency & The unit of privacy / what changes between worlds \\
Budget type & Odds inflation / max-divergence / TV / tradeoff curve \\
Per-epoch budget & $(\eta,\delta)$ or equivalent \\
Threat/window + usage & \texttt{tw\_id} plus the usage/horizon assumption and accounting regime (rulebook: Synthesis~19) \\
Owned proofs & Which papers/certificates provide the proof objects \\
Artifact digests & Hash pointers to evidence (on request) \\
\bottomrule
\end{tabular}
\caption{OINL is a derived label: it summarizes (and links to) ABOM fields.}
\end{table}

\paragraph{OINL and ``referee-grade brevity.''}
The papers in this archive can stay short because OINL externalizes the long tail of artifacts.
A referee can demand the archive; a casual reader can still interpret the claim.

\section{Interfacing with receipts, certificates, and transparency}

Receipt line items (Synthesis~12) are the publication-level companion of ABOM/OINL: each receipt line item repeats the same claim fields (\texttt{exposure\_nf\_id}, \texttt{tw\_id}, and when applicable \texttt{plan\_spec\_id} and \texttt{state\_decl\_id}) so that a reader can compare the OINL against the receipt without opening private artifacts.\cite{series:receiptitems}

\paragraph{Horizon claims are structured.}
If an ABOM/OINL entry quotes a horizon-$n$ total ("per day", "for $Q$ lookups"), it should also name the usage assumption and bind the accountant via \texttt{plan\_spec\_id} (and \texttt{state\_decl\_id} when state enters), following the receipt-accounting rulebook.\cite{series:accountingrulebook}


\paragraph{A drift rule worth making explicit.}
If \texttt{plan\_spec\_id} or \texttt{state\_decl\_id} changes, the publication must mint a new receipt/ABOM entry that pins the new digest-bound semantics.
Whether the underlying privacy claim needs recertification depends on whether the change is merely ``publication-only'' (re-run the verifier) or claim-affecting; we centralize the taxonomy in Synthesis~18 so other papers do not restate it.\cite{series:planstateequiv}

ABOM is compatible with (and intentionally defers details to) the series' receipt/certificate pipeline:
\begin{itemize}[leftmargin=*]
\item \textbf{Release receipts:} Release~A defines a ledger bundle and a compact PVSA summary attestation.\cite{series:releaseproperty}
ABOM can point to the bundle and include the PVSA digest.
\item \textbf{Congestion certificates:} W-Congestion-EQ defines congestion-specific certificate fields and log anchoring patterns.\cite{series:wcongeq}
ABOM can embed the certificate digest and an inclusion-proof digest.
\item \textbf{Optimization certificates:} BossFight~C and the 2026-01-29 wiki post treat dual certificates as proof objects for TV privacy designs.
ABOM can include the primal/dual transcript digests.
\end{itemize}

\section{Takeaways}
\begin{itemize}[leftmargin=*]
\item \textbf{ABOM is the commitment:} a signed manifest that pins (by digest) the artifacts and the \emph{semantics} of the claim (projection, adjacency, budget representation, horizon), plus any replay and state surfaces the guarantee depends on.
\item \textbf{OINL is the diffable summary:} a compact label extracted from ABOM so reviewers and downstream users can compare claims across releases without unpacking the full artifact archive.
\item \textbf{Source-only can still be referee-grade:} publish ABOM/OINL in the paper; ship the heavy archive on request; let third parties verify the delivered bundle matches the published digests.
\end{itemize}

\section{Limitations and scope}
ABOM/OINL do not prevent lying: a publisher could sign a manifest for a claim that is false.
Their purpose is narrower: make claims \emph{auditable} and \emph{comparable} across releases, and make it straightforward for others to request and verify the full artifact archive.

\begin{thebibliography}{99}\small

\bibitem{series:artifactpolicy}
Anonymous.
\newblock Artifact and Disclosure Policy for the One-True-Archive (Source-Only Public Release).
\newblock Series synthesis note (Synthesis~6), 2026. (This archive.)

\bibitem{series:oddsinflation}
Anonymous.
\newblock Odds Inflation, Privacy Loss, and Endpoint Anonymity Bounds.
\newblock Anonymity series Paper~A, 2026. (This archive.)

\bibitem{series:endpointbridge}
Anonymous.
\newblock Endpoint Metrics Bridge: Odds Inflation, PML/MaxL, and Identification Sizing.
\newblock Series synthesis note (Synthesis~5), 2026. (This archive.)

\bibitem{series:releaseproperty}
Anonymous.
\newblock Privacy as a Release Property: Proof-Carrying Anonymity Receipts for Anonymous DHTs.
\newblock Release and destination Paper~A, 2026. (This archive.)

\bibitem{series:wcongeq}
Anonymous.
\newblock W-Congestion-EQ: Proof-Carrying Congestion Budgets with Certificates and Transparency Logs for Anonymous Overlays.
\newblock Congestion series Paper~3, 2026. (This archive.)

\bibitem{series:traceifaces}
Anonymous.
\newblock Trace Interfaces for Anonymous DHTs: Observation Tiers, Committee-Contact Traces, and Exposure Normal Forms.
\newblock Series synthesis note (Synthesis~3), 2026. (This archive.)

\bibitem{series:threatwindows}
Anonymous.
\newblock Threat Models and Repetition Windows for Anonymous DHTs: A Short Declaration Checklist for Referee-Grade Budget Claims.
\newblock Series synthesis note (Synthesis~4), 2026. (This archive.)

\bibitem{series:accountingrulebook}
Anonymous.
\newblock Receipt Accounting and Horizon Composition for Anonymous-DHT Budgets.
\newblock Series synthesis note (Synthesis~19), 2026. (This archive.)

\bibitem{series:contractobjects}
Anonymous.
\newblock Contract Objects for Anonymous-DHT Privacy Budgets and Claims.
\newblock Series synthesis note (Synthesis~0), 2026. (This archive.)

\bibitem{series:receiptitems}
Anonymous.
\newblock Receipt Line Items for Anonymity Claims: A Minimal Tuple Schema for Published Budgets (Series Synthesis).
\newblock Series synthesis note (Synthesis~12), 2026. (This archive.)

\bibitem{series:planstateequiv}
Anonymous.
\newblock Digest-Bound Replay Plans and State Declarations: Equivalence, Minimality, and Change Control.
\newblock Series synthesis note (Synthesis~18), 2026. (This archive.)

\bibitem{series:state1}
Anonymous.
\newblock State-Dependent Anonymity: Selection and State Evolution as an Attack Surface.
\newblock AnonDHT State series Paper~1, 2026. (This archive.)

\bibitem{series:state3}
Anonymous.
\newblock Closed-View Auditing for Stateful Anonymity: Odometers, Alarms, and Machine-Checkable CPPCs.
\newblock AnonDHT State series Paper~3, 2026. (This archive.)

\bibitem{ntiaSBOM2019}
National Telecommunications and Information Administration (NTIA).
\newblock Survey of Existing SBOM Formats and Standards.
\newblock Whitepaper, Oct.~2019.\;
\newblock \url{https://www.ntia.gov/sites/default/files/publications/ntia_sbom_formats_and_standards_whitepaper_-_version_20191025_0.pdf}.

\bibitem{cyclonedx}
OWASP Foundation.
\newblock CycloneDX Bill of Materials (BOM) Standard.
\newblock \url{https://cyclonedx.org/}. (Accessed Feb.~2026.)

\bibitem{spdx301}
The Linux Foundation (SPDX Project).
\newblock SPDX Specification v3.0.1.
\newblock 2024.\;\url{https://spdx.dev/wp-content/uploads/sites/31/2024/12/SPDX-3.0.1-1.pdf}.

\bibitem{intoto2019}
S.~Torres-Arias, H.~Afzali, T.~K.~Kuppusamy, R.~Curtmola, and J.~Cappos.
\newblock In-Toto: Providing farm-to-table guarantees for bits and bytes.
\newblock In \emph{USENIX Security}, 2019, pp.~1393--1410.\;
\newblock \url{https://www.usenix.org/conference/usenixsecurity19/presentation/torres-arias}.
\end{thebibliography}

\end{document}
